%-------------------------
% Resume - Siddardth Pathipaka
% Version B: Process / Continuous Improvement Engineer
%-------------------------
\documentclass[letterpaper,9pt]{extarticle}

\usepackage[utf8]{inputenc}
\usepackage[T1]{fontenc}
\usepackage{geometry}
\usepackage{enumitem}
\usepackage{titlesec}
\usepackage{hyperref}
\usepackage{xcolor}
\usepackage{fontspec}
\usepackage{tabularx}

\geometry{
  letterpaper,
  left=0.47in,
  right=0.47in,
  top=0.20in,
  bottom=0.15in,
  noheadfoot
}

\setmainfont{Carlito}[
  BoldFont=Carlito Bold,
  ItalicFont=Carlito Italic,
  BoldItalicFont=Carlito Bold Italic
]

\pagestyle{empty}

\hypersetup{
  colorlinks=true,
  linkcolor={rgb:red,5;green,99;blue,193},
  urlcolor={rgb:red,5;green,99;blue,193},
}
\definecolor{linkblue}{RGB}{5,99,193}

\titleformat{\section}{
  \bfseries\fontsize{10pt}{12pt}\selectfont
}{}{0em}{}[\vspace{-4pt}\rule{\textwidth}{0.5pt}]
\titlespacing*{\section}{0pt}{4pt}{1pt}

\setlist[itemize]{
  leftmargin=0.25in,
  labelsep=0.06in,
  itemsep=0pt,
  parsep=0pt,
  topsep=0pt,
  label=\textbullet
}

\newcommand{\expentry}[2]{%
  \vspace{2pt}%
  \noindent\begin{tabularx}{\textwidth}{@{}X r@{}}
    {\fontsize{11pt}{13pt}\selectfont\textbf{#1}} & {\fontsize{11pt}{13pt}\selectfont\textbf{#2}}
  \end{tabularx}%
  \vspace{-1pt}%
}

\newcommand{\eduentry}[2]{%
  \vspace{2pt}%
  \noindent\begin{tabularx}{\textwidth}{@{}X r@{}}
    {\fontsize{11pt}{13pt}\selectfont\textbf{#1}} & {\fontsize{11pt}{13pt}\selectfont\textbf{#2}}
  \end{tabularx}%
}

\newcommand{\skillline}[2]{%
  \noindent{\fontsize{9pt}{11pt}\selectfont\textbf{#1} #2}\par\vspace{1pt}%
}

\begin{document}

%----------NAME----------
\begin{center}
  {\fontsize{14pt}{16pt}\selectfont\textbf{Siddardth Pathipaka}}\\[2pt]
  {\fontsize{10pt}{12pt}\selectfont\textbf{siddardth1524@gmail.com \textbar\ +1(217) 255-0104 \textbar\ %
  \href{https://github.com/siddardth7}{\color{linkblue}\underline{GitHub}} \textbar\ %
  \href{https://linkedin.com/in/siddardth-pathipaka}{\color{linkblue}\underline{LinkedIn}} \textbar\ %
  \href{https://siddardth7.github.io/portfolio/}{\color{linkblue}\underline{Portfolio}}}}
\end{center}

\vspace{-4pt}

%----------SUMMARY----------
\section{SUMMARY}
{\fontsize{9pt}{11pt}\selectfont
\textbf{%%SUMMARY%%}
}

%----------SKILLS----------
\section{SKILLS}
\vspace{1pt}
%%SKILLS_BLOCK_START%%
\skillline{Continuous Improvement:}{FMEA, SPC, 8D Root Cause Analysis, CAPA, Lean Principles, Defect Reduction, Process Repeatability, Yield Improvement}
\skillline{Process Engineering:}{Cure Cycle Development, Assembly Optimization, Workflow Redesign, First Article Inspection, MRB Disposition, GD\&T, CMM Inspection}
\skillline{Composite Processing:}{Prepreg Layup, Autoclave Processing, Vacuum Bagging, Debulk, Out-of-Autoclave Methods}
\skillline{Simulation \& Software:}{ABAQUS, ANSYS Fluent, FEA, SolidWorks, MATLAB, Python, AutoCAD}
%%SKILLS_BLOCK_END%%

%----------EXPERIENCE----------
\section{MANUFACTURING \& ENGINEERING EXPERIENCE}

% SAMPE -- lead card for CI
\expentry{Manufacturing Lead \textbar\ SAMPE Composite Fuselage Program, UIUC}{Jan 2025 \textbar\ May 2025}
\begin{itemize}
  \item Implemented \textbf{FMEA} on composite cure process, identifying vacuum bag leaks as highest-risk failure mode \textbf{(RPN=60)} and deploying automated leak detection -- achieving \textbf{zero process deviations} across all cure cycles.
  \item Redesigned layup workflow by switching to a \textbf{continuous prepreg blanket} approach, eliminating delamination risk and establishing \textbf{repeatable process standards} for the full production sequence.
  \item Used \textbf{simulation-informed parameter optimization} (ABAQUS FEA + Classical Lamination Theory) to select a stacking sequence that reduced deflection by \textbf{38\%} at \textbf{150\% design load (1,000 lbf)} -- replacing a trial-and-error build cycle with a single validated configuration.
  \item Delivered first-article composite fuselage with \textbf{under 2\% void content} using \textbf{autoclave cure (275°F, 40 psi)}, meeting all structural acceptance criteria on first build.
\end{itemize}

% Tata Boeing
\expentry{Manufacturing \& Quality Intern \textbar\ Tata Boeing Aerospace}{Oct 2022 \textbar\ Mar 2023}
\begin{itemize}
  \item Ran \textbf{SPC-based process investigation} isolating \textbf{CNC tool wear} as root cause of position tolerance failures, implementing corrective tooling change intervals that reduced defect rate from \textbf{15\% to under 3\%} -- with the fix adopted into standing production controls.
  \item Drove cross-functional \textbf{8D root cause analysis} on GE engine component nonconformances, cutting MRB cycle time by \textbf{22\%} and improving disposition turnaround.
  \item Implemented \textbf{FMEA} on supplier nonconformance that prevented \textbf{\textasciitilde\$3,000} in scrap through engineering-approved use-as-is disposition, reducing material waste.
  \item Standardized \textbf{CMM inspection workflow} for \textbf{450+} flight-critical components to \textbf{0.02 mm accuracy}, achieving \textbf{zero customer escapes} on GE and Boeing programs.
\end{itemize}

% Graduate Research
\expentry{Graduate Research Assistant \textbar\ Beckman Institute, UIUC}{May 2024 \textbar\ Dec 2024}
\begin{itemize}
  \item Developed a \textbf{rapid composite cure process} using frontal polymerization, reducing processing cycle time from \textbf{8+ hours to under 5 minutes} -- an \textbf{out-of-autoclave} approach that demonstrates a path to eliminating autoclave-constrained production cycles.
  \item Built computational process models predicting cure behavior within \textbf{10\% velocity accuracy} and \textbf{3°C} of measured peak temperatures, compressing process parameter optimization cycle by \textbf{94\%}.
\end{itemize}

% EQIC
\expentry{Manufacturing Intern \textbar\ EQIC Dies \& Moulds Engineers Pvt. Ltd.}{Jun 2022 \textbar\ Aug 2022}
\begin{itemize}
  \item Mapped \textbf{12-stage die production workflow} -- from CNC machining through EDM, heat treatment, and final assembly -- identifying inter-stage handoff points as primary sources of tolerance accumulation and quality risk.
  \item Verified die tolerances to \textbf{±0.02 mm} against \textbf{GD\&T specs} and validated \textbf{parting surface alignment ({>}80\% contact)} to prevent downstream casting defects in high-pressure tooling.
\end{itemize}

%----------PROJECTS----------
\section{PROJECTS}

\expentry{Brakes \& Suspension Lead \textbar\ Electric Solar Vehicle Championship (ESVC)}{Oct 2019 \textbar\ Apr 2023}
\begin{itemize}
  \item Led design, manufacturing, and assembly of brake and suspension hardware across a \textbf{30-member} cross-functional team, achieving \textbf{zero mechanical failures} in competition.
  \item Validated process output with \textbf{180.43/200 on Cross Pad Test} (top handling score), demonstrating manufactured hardware performance.
\end{itemize}

\expentry{Structures \& Manufacturing Lead \textbar\ SAE Aero Design Challenge}{Dec 2021 \textbar\ Sep 2022}
\begin{itemize}
  \item Developed and executed \textbf{assembly sequencing plan} for a \textbf{92-inch tapered balsa wing}, coordinating reinforcement integration and \textbf{CG balancing at 25\% MAC} across multiple build phases.
\end{itemize}

\expentry{Research Assistant \textbar\ VNR VJIET Composites Lab, Hyderabad}{Aug 2022 \textbar\ Apr 2023}
\begin{itemize}
  \item Optimized composite formulation across \textbf{24 experimental variants} using powder metallurgy, achieving \textbf{27\% hardness improvement (548 to 699 HV)} through systematic reinforcement parameter testing.
\end{itemize}

%----------EDUCATION----------
\section{EDUCATION}

\eduentry{M.S., Aerospace Engineering \textbar\ University of Illinois at Urbana-Champaign}{Jan 2024 \textbar\ Dec 2025}
\vspace{1pt}
\hspace{0.12in}{\fontsize{9pt}{11pt}\selectfont\textbf{\textit{Relevant Coursework:}} \textit{Manufacturing Processes, Composite Materials, Materials Characterization, Structural Mechanics, Finite Element Methods}}

\eduentry{B.Tech, Mechanical Engineering \textbar\ VNR VJIET}{Aug 2019 \textbar\ May 2023}

\end{document}
